\section{Conclusão}
% Neste trabalho investigamos a calibração de sensores de baixo custo para estimação da concentração horária de NO$_2$ a partir do conjunto de dados Air Quality UCI. Após pré-processamento com tratamento de valores ausentes, transformação logarítmica de variáveis assimétricas e padronização Z-score, avaliamos quatro abordagens: regressão linear ordinária (OLS), regressão Ridge, regressão por mínimos quadrados parciais (PLS) e uma rede neural MLP simples.

% Os resultados apresentados na Seção~\ref{sec:resultados} mostram que os modelos lineares regularizados alcançam desempenho elevado e consistente. Tanto Ridge quanto PLS obtiveram RMSE de teste em torno de 10{,}4 e coeficiente de determinação $R^2$ próximo de 0{,}89, superando ligeiramente a OLS em estabilidade e erro médio, ao mesmo tempo em que mantêm interpretabilidade e baixo custo computacional. A rede neural MLP apresentou desempenho inferior, com maior erro e variabilidade entre folds, o que indica que a complexidade não linear não se justifica para o tamanho de amostra e configuração adotados.\newline

% Do ponto de vista de interpretação, a análise conjunta dos coeficientes padronizados do Ridge e dos componentes do PLS indica que um subconjunto de variáveis --- em particular sensores fortemente relacionados a NOx/NO$_2$ e algumas variáveis meteorológicas --- responde pela maior parte da capacidade preditiva. Esse achado é coerente com a literatura sobre calibração de sensores de baixo custo e reforça a utilidade de métodos que lidam explicitamente com multicolinearidade, seja via penalização (Ridge) ou via componentes supervisionados (PLS).\newline

% Entre as principais limitações destacam-se o uso de dados de uma única estação em um período específico (2004--2005), o que restringe a generalização para outros contextos, e a ausência de exploração sistemática de arquiteturas neurais mais profundas ou modelos temporais explícitos. Como perspectivas futuras, destacam-se a avaliação de modelos multialvo que calibrem simultaneamente diferentes poluentes, a incorporação de informação temporal mais rica em modelos lineares e não lineares e a validação em cenários reais de implantação de redes de sensores, incluindo deriva temporal e mudanças no padrão de emissões.\newline

% Em síntese, os resultados indicam que modelos lineares bem regularizados, combinados a um pré-processamento adequado, constituem uma solução robusta e eficiente para calibração de sensores de baixo custo de NO$_2$, oferecendo um compromisso atraente entre acurácia, complexidade e interpretabilidade em aplicações de monitoramento da qualidade do ar.
